% !TEX program = xelatex
\documentclass[11pt,a4paper]{article}

\usepackage{fontspec}
\setmainfont{TeX Gyre Pagella}
\usepackage{microtype}
\usepackage{amsmath,amsthm,amssymb}
\usepackage{geometry}
\geometry{margin=1.25in}
\usepackage{setspace}
\onehalfspacing
\usepackage{hyperref}
\hypersetup{colorlinks=true, linkcolor=black, citecolor=black, urlcolor=black}
\usepackage{csquotes}
\usepackage{parskip}
\setlength{\parskip}{0.6em}
\setlength{\parindent}{1.5em}
\emergencystretch=2em

\theoremstyle{plain}
\newtheorem{theorem}{Theorem}[section]
\newtheorem{proposition}[theorem]{Proposition}
\newtheorem{lemma}[theorem]{Lemma}
\newtheorem{corollary}[theorem]{Corollary}

\theoremstyle{definition}
\newtheorem{definition}[theorem]{Definition}
\newtheorem{remark}[theorem]{Remark}

\title{\textbf{Propagating Constraint Fields and Dynamic Structural Redistribution:\\
\large Toward a Field-Theoretic Interpretation of Neural\\
Heterogeneity and Collective Cell Migration}}
\author{Flyxion\\[0.5em]
\textit{Independent Researcher}}
\date{May 2026}

\begin{document}

\maketitle

\begin{abstract}
Two recently published empirical studies, one investigating structural brain aging in
generalized anxiety disorder and the other characterizing the dynamic redistribution of
the tight-junction protein ZO-1 during collective epithelial cell migration, exhibit a
shared organizational structure that resists explanation by standard localist or
entity-centered biological models. In the brain-aging study, individuals with
generalized anxiety disorder do not uniformly display the accelerated aging signature
that disease-model reasoning would predict; instead they exhibit dramatically increased
variance in predicted age difference, suggesting that the disorder broadens the space of
accessible neural trajectories rather than displacing the entire population along a single
developmental axis. In the cell-migration study, ZO-1 does not remain fixed at
apical junctional complexes but dynamically shuttles to podosome-like basal structures
in response to propagating ERK activation waves, reorganizing the tissue's mechanical
and signaling architecture collectively across space and time. We argue that both
phenomena are most coherently interpreted within a generalized field-theoretic framework
grounded in the Relativistic Scalar-Vector Plenum formalism, in which local organizational
states are characterized by a scalar coherence field, a vector transport field, and an
entropic accessibility measure. Under this interpretation, psychiatric heterogeneity
becomes a widening of admissible trajectory distributions rather than a categorical
pathological shift, while collective tissue migration becomes a case of constraint
redistribution driven by propagating activation fronts. The analysis yields specific
structural predictions about variance geometry, wave-coupled phase transitions between
stabilizing and exploratory organizational regimes, and the general insufficiency of
mean-field biomarker approaches for systems governed by distributed constraint dynamics.
\end{abstract}

\tableofcontents
\newpage

\section{Introduction: From Static Biomarkers to Dynamical Constraint Fields}

A strongly localist approach to biological explanation treats macroscopic phenomena as
the aggregated result of discrete components interacting according to intrinsic stable
properties. This entity-centered framework has produced genuine successes, particularly
in structural biochemistry and classical genetics, but it encounters systematic difficulty
whenever the explanatory target is coherent large-scale behavior that emerges from
distributed interaction rather than from the characteristics of any single component.
Two recent empirical contributions, one in translational psychiatry and one in cell
biology, expose precisely this limitation while simultaneously pointing toward an
alternative organizing principle. A third contribution, in computational neuroscience,
provides what may be the most direct empirical instantiation of the organizational
geometry the others implicitly suggest.

The first paper, by Richier and collaborators within the ENIGMA-Anxiety consortium,
investigates whether generalized anxiety disorder (GAD) is associated with accelerated
structural brain aging as measured by convolutional neural network estimation of
chronological age from multimodal MRI data~\cite{Richier2026}. The core metric is the
predicted age difference (PAD), computed as estimated brain age minus actual chronological
age, with positive values indicating an apparently older structural profile. The central
result contradicts the original hypothesis: individuals with GAD do not show significantly
elevated mean PAD relative to healthy controls. What they do show, particularly among
adults older than twenty-five, is substantially increased \emph{variance} in PAD.
Some individuals with GAD display neural profiles consistent with accelerated aging;
others display profiles consistent with delayed or atypical development. The diagnostic
category does not map to a single point in developmental space but rather to a broadened
and heterogeneous cloud.

The second paper, by Hirano and collaborators, examines how the tight-junction scaffolding
protein ZO-1 participates in collective epithelial migration~\cite{Hirano2026}. Using
live-cell imaging of migrating MDCK epithelial monolayers, the authors demonstrate that
ZO-1 is not a static structural component of cell-cell junctions but a dynamic participant
in tissue-scale coordination. During migration, ERK activation propagates across the
epithelial sheet as traveling waves; in their wake, ZO-1 detaches from apical junctional
complexes and accumulates at basal podosome-like structures where it influences traction
force generation and extracellular matrix remodeling. These redistributions themselves
propagate collectively. The tissue does not consist of cells individually deciding to
migrate; it behaves as a coupled dynamical medium through which activation fronts
reorganize mechanical architecture across space and time.

The third paper, by Wang, Lou, Wei, and collaborators, introduces a geometry-aware
framework for EEG/MEG source imaging that reconstructs whole-brain neural dynamics
by decomposing the individual cortical surface into Laplace--Beltrami eigenmodes,
termed geometric basis functions (GBFs)~\cite{Wang2026}. The central empirical finding
is that whole-brain activity observed across approximately 32,000 cortical vertices in
fMRI can be reconstructed to high fidelity using only 200--300 such modes. Neural
sources are expressed as $x(t) = \sum_i \theta_i(t)\psi_i$, where $\psi_i$ are
geometry-derived eigenfunctions rather than arbitrary basis vectors. The inverse
solution is regularized by the spectral prior
$\Sigma^{-1} = \mathrm{diag}(-\beta/\log\lambda_i)$, which penalizes high-frequency
modes increasingly, thereby embedding an explicit admissibility weighting into the
reconstruction operator. The paper further tracks stimulation-evoked propagating cortical
waves using optical-flow analysis, recovering directed propagation streamlines and
convergence hubs consistent with known anatomical pathways. This framework operates
directly on a Riemannian manifold (the cortical surface) with intrinsic differential
operators, making it the empirical system in this essay most closely aligned with the
abstract RSVP field equations.

The structural parallel across these three findings is not superficial. In each case,
an entity-centered account that assigns stable properties to diagnosable individuals,
to protein molecules, or to individual cortical sources fails to capture the dominant
organizational signal. In each case, the explanatory work is done by distributed
field-like dynamics: trajectory-space variance geometry in the psychiatric case;
traveling wave-coupled phase transitions in the cell biology case; and geometry-constrained
eigenmode compression with wave propagation in the neuroimaging case. The present essay
develops these parallels systematically, arguing that all three phenomena can be
interpreted within a common generalized field-theoretic framework without reducing
any to the others or making implausible claims about ontological identity across
biological scales.

The interpretive framework we employ, the Relativistic Scalar-Vector Plenum (RSVP)
formalism, was developed originally in the context of cosmological field theory and has
been extended progressively to problems in cognition, political economy, and collective
computation. We do not claim that the cellular ERK system, the anxious brain, and the
cortical eigenmode structure are literally implementing the same physical equations.
We claim something more modest and more defensible: that the organizational principles
visible in all three systems are instances of a common mathematical structure, namely
the propagation of constraint redistribution across coupled fields defined over
organizational manifolds, and that formalizing this structure yields genuine predictive
and explanatory leverage that entity-centered models cannot provide.

Section~\ref{sec:gad} develops the RSVP-theoretic interpretation of the GAD brain-age
findings. Section~\ref{sec:erk} analyzes the ZO-1 and ERK wave dynamics.
Section~\ref{sec:phase} develops a unified account of structural phase transitions between
stabilizing and exploratory organizational regimes. Section~\ref{sec:admissibility}
establishes admissibility as a cross-domain mathematical primitive. Section~\ref{sec:rsvp}
presents the formal framework including interpretation of all three empirical systems.
Section~\ref{sec:implications} draws implications for psychiatry, developmental biology,
and collective intelligence research. Section~\ref{sec:clarifications} addresses
conceptual boundaries, theoretical distinctions, and the scope of the framework's claims.
Section~\ref{sec:limitations} acknowledges limitations explicitly.
Section~\ref{sec:predictions} derives experimentally discriminable predictions.
Section~\ref{sec:conclusion} concludes.

\section{Brain-Age Variability and Heterogeneous Neural Trajectories}
\label{sec:gad}

\subsection{The PAD Geometry of Anxiety}

The brain-age framework rests on a simple and elegant idea: a machine learning model
trained to estimate chronological age from structural neuroimaging can be applied to
clinical populations, and the residual discrepancy between predicted and actual age
provides a measure of how structurally atypical that population's brains are relative
to healthy developmental norms~\cite{Richier2026}. Positive predicted age difference
implies a neurologically older-appearing structure; negative PAD implies relative youth
or atypicality in a different direction. The implicit model is that pathology acts as an
accelerant, pushing brains prematurely toward configurations typical of older healthy
individuals.

This model is intuitively appealing and has shown positive results in conditions such as
schizophrenia and major depression. Its failure in GAD is therefore informative. The
absence of a significant mean PAD elevation does not merely show that GAD is neurologically
undetectable at the group level; it reveals that the appropriate description of the
GAD population in trajectory space is not a translation (a shift of the mean) but a
dilation (an expansion of the variance). The population is more spread out in the
space of structural brain configurations, not more concentrated in a pathological direction.

Let $\mathcal{T}$ denote the space of structural brain configurations reachable across
a healthy lifespan, parameterized by chronological age $t \in [t_0, t_\infty]$ and
by a high-dimensional morphometric state vector $\mathbf{x} \in \mathbb{R}^n$.
The healthy control population traces a distribution $\mu_{HC}$ over $\mathcal{T}$,
which the CNN model approximates by learning to invert the projection
$\pi: \mathcal{T} \to [t_0, t_\infty]$ that extracts chronological age from structural
configuration. The PAD for an individual is then $\hat{t}(\mathbf{x}) - t$,
where $\hat{t}$ is the learned inverse map.

The Richier et al.\ findings can be formalized as follows.

\begin{definition}[PAD Geometry]
Let $\mu_{HC}$ and $\mu_{GAD}$ denote the distributions of structural configurations
in healthy controls and GAD individuals respectively, projected to PAD values via
$\delta(t, \mathbf{x}) = \hat{t}(\mathbf{x}) - t$. The GAD population exhibits
\emph{trajectory-space dilation} rather than translation if
\[
\mathbb{E}_{\mu_{GAD}}[\delta] \approx \mathbb{E}_{\mu_{HC}}[\delta]
\quad \text{while} \quad
\mathrm{Var}_{\mu_{GAD}}[\delta] > \mathrm{Var}_{\mu_{HC}}[\delta],
\]
with the inequality becoming more pronounced in adult subpopulations.
\end{definition}

This is precisely what Richier et al.\ report. The finding suggests that GAD does not
primarily operate by shifting the brain's developmental trajectory in a single direction
but by increasing the dispersion of trajectories across individuals. Some individuals
consolidate neural architecture prematurely in response to chronic threat-anticipation;
others remain in more developmentally diffuse or plastic states. The disorder does not
produce a single neurological signature because it does not impose a single developmental
pressure but rather amplifies the sensitivity of developmental trajectories to individual
variation in symptom severity, comorbidity, medication, and life history.

\subsection{Admissibility Geometry and Trajectory Variance}

The entity-centered biomarker model implicitly assumes that GAD is a fixed point in
diagnostic space corresponding to a fixed region in brain-configuration space. The PAD
results falsify this assumption. A more adequate model treats the condition as a
deformation of the admissibility geometry of developmental trajectories: a change in
which configurations are dynamically accessible, and with what probability, rather than
a change in the mean configuration.

\begin{definition}[Admissibility Field]
An admissibility field over $\mathcal{T}$ is a non-negative function
$S: \mathcal{T} \to [0, \infty)$ that assigns to each configuration $\mathbf{x}$ at age
$t$ the local volume of accessible future trajectories. High $S$ corresponds to high
developmental plasticity or instability; low $S$ corresponds to a consolidated, constrained
future trajectory.
\end{definition}

In healthy development, $S$ typically decreases with age as neural architecture
consolidates. Critical periods are intervals during which $S$ is locally elevated,
permitting major structural reorganization. Pathological processes can alter $S$ in two
ways: by globally elevating it (increasing plasticity or instability), by globally
suppressing it (premature consolidation), or by increasing its \emph{variance} across
individuals without strongly shifting its mean.

The third pattern is what the GAD results suggest. The disorder does not uniformly
elevate or suppress developmental plasticity; it increases the heterogeneity of
admissibility profiles across individuals. Some GAD brains show signatures of suppressed
$S$ (premature consolidation, positive PAD); others show elevated $S$ (delayed
consolidation, negative PAD); the population-level mean remains close to zero.

\begin{proposition}
If GAD acts primarily through increased variance in the admissibility field
rather than through a directional shift in its mean, then the PAD distribution of
the GAD population will be approximately mean-zero but with elevated variance relative
to controls, and symptom severity will correlate with the magnitude of individual
PAD rather than with its sign.
\end{proposition}

\begin{proof}[Argument]
Let $S_i$ denote the admissibility profile of individual $i$. If GAD shifts $S_i$
by a zero-mean random perturbation $\xi_i$ with $\mathrm{Var}[\xi_i] > \mathrm{Var}[\xi^{HC}]$,
then the resulting PAD distribution inherits the increased variance without inheriting
a directional mean shift. Symptom severity, as a monotone function of the perturbation
magnitude $|\xi_i|$, will then correlate with $|\delta_i|$ regardless of sign.
The exploratory analyses reported by Richier et al.\ showing symptom-PAD correlations
after controlling for medication and comorbidity are consistent with this structure.
\end{proof}

\subsection{The Insufficiency of Univariate Reduction}

A secondary finding of Richier et al.\ is that their CNN-based multivariate approach
detects the heterogeneity signal more clearly than traditional univariate neuroimaging
analyses would. This is not a minor methodological footnote. It reflects a deep point
about the relationship between representational dimensionality and the detectability of
distributed organizational structure.

Traditional univariate analyses project the high-dimensional brain-con\-fig\-u\-ra\-tion manifold
$\mathcal{T}$ to a collection of scalar measurements at individual voxels or regions.
These projections preserve local amplitude information but destroy the distributed
relational structure that encodes developmental state. The CNN model, by learning
spatially distributed nonlinear functions of the full three-dimensional input, retains
more of this relational structure and therefore detects organizational patterns invisible
to univariate methods.

This is exactly what should be expected if the relevant organizational signal is carried
by the admissibility geometry of the manifold rather than by the marginal distributions
at individual measurement sites. Variance in trajectory space is a property of the
manifold's global structure; collapsing the manifold to independent scalars destroys
precisely the information that makes this variance detectable.

One important epistemic limitation of the trajectory-space dilation interpretation
deserves explicit acknowledgment. The Richier et al.\ data are cross-sectional: each
individual contributes a single PAD measurement at a single time point. The dilation
interpretation treats increased cross-sectional PAD variance as evidence for increased
variance in developmental trajectories, but this inference requires the additional
assumption that the cross-sectional distribution reflects the underlying longitudinal
dynamics. Individuals with high positive PAD may have followed a consistently
accelerated trajectory, or may have undergone a recent acute consolidation event; the
cross-sectional measurement cannot distinguish these. Longitudinal imaging designs,
in which the same individuals are followed across developmental transitions, would be
required to confirm that the trajectory-space dilation is a genuine dynamical feature
rather than an artifact of population heterogeneity in initial developmental state.
This limitation motivates the experimental predictions developed in Section~\ref{sec:predictions}.

\section{ERK Activation Waves as Distributed Coordination Dynamics}
\label{sec:erk}

\subsection{The Tissue as a Coupled Dynamical Medium}

The standard model of epithelial cell migration treats cells as semi-autonomous agents
that respond to local chemical gradients, express motility machinery, and collectively
produce directed movement as a statistical aggregate of individual decisions. This model
works reasonably well for loosely coupled populations but becomes inadequate when
describing the highly coordinated collective migration observed in wound healing,
embryonic development, and cancer invasion. The Hirano et al.\ paper on ZO-1 and ERK
wave dynamics provides a particularly clear demonstration of why~\cite{Hirano2026}.

The central observation is that ERK activation does not arise independently in each cell
in response to local stimuli but propagates across the epithelial monolayer as a traveling
wave. The tissue, as a coupled dynamical medium, supports the propagation of biochemical
activation fronts in ways that individual cells, considered in isolation, cannot generate.
This propagating activation then triggers a second-order dynamical event: the redistribution
of ZO-1 from apical junctional complexes to basal podosome-like structures. This
redistribution itself propagates collectively, following the ERK wave with a characteristic
delay.

The result is a tissue that reorganizes its structural architecture in a wave-like fashion
across space and time. The tight-junction protein that normally maintains epithelial
barrier integrity transiently relocates to structures associated with invasive migration
and extracellular matrix remodeling. The collective consequence is directed migration
driven not by independent cellular decisions but by a distributed signaling computation
performed across the entire monolayer.

\subsection{ERK Waves as Constraint-Propagating Fronts}

From a field-theoretic perspective, the ERK activation wave is best understood not as a
simple diffusion front but as a constraint-propagating front: a traveling region of local
dynamical reorganization that alters the admissibility structure of each cell it passes
through, enabling configurations (podosome formation, ZO-1 redistribution, elevated
traction force) that were not accessible before the wave arrived.

\begin{definition}[Constraint-Propagating Front]
A constraint-propagating front in a coupled spatial dynamical system is a traveling
region $\mathcal{F}(t) \subset \Omega$ such that passage of $\mathcal{F}$ through a
local neighborhood $U$ alters the local admissibility structure $S|_U$, enabling
transitions between organizational regimes that are suppressed in the pre-wave state.
\end{definition}

The ERK wave is a constraint-propagating front because it enables the ZO-1 redistribution
transition: MEK inhibition, which suppresses ERK wave propagation, also suppresses ZO-1
translocation to podosomes, demonstrating that the wave is not merely correlated with
the structural reorganization but is causally required for it~\cite{Hirano2026}.
The wave opens a dynamical window during which the cell can exit the tight-junction
maintenance regime and enter the podosome-mediated invasive regime.

This is structurally different from simple diffusion or gradient-driven migration.
In diffusion, the signal spreads to enable a quantitative response whose direction is
determined by concentration. In constraint propagation, the wave alters the qualitative
topology of local accessible state space: it enables regime transitions rather than
merely modulating continuous response intensity. The analysis of admissible and
non-classical waves in hyperbolic systems provides formal precedent for precisely this
kind of topology-changing wave structure, in which the wave selects among formally
permissible branches of evolution rather than simply transporting energy~\cite{ElAbbassi2025}.

\subsection{Wave Mechanics and Propagation Velocity}

The mathematical structure of ERK waves in epithelial systems has features consistent
with reaction-diffusion dynamics of the excitable medium type~\cite{Murray2003}. A
minimal model treats local ERK activation $u(\mathbf{r}, t)$ as satisfying

\begin{equation}
\frac{\partial u}{\partial t} = D \nabla^2 u + f(u, v),
\quad
\frac{\partial v}{\partial t} = \epsilon g(u, v),
\label{eq:erk_wave}
\end{equation}

where $v$ is a recovery variable, $f$ and $g$ are nonlinear kinetic functions
specifying the local excitable dynamics, $D$ is the effective coupling coefficient
(which in an epithelial sheet combines intracellular signaling diffusion with
paracrine coupling through gap junctions and extracellular ligands), and $\epsilon \ll 1$
is the timescale separation between the fast activation variable and the slow recovery.
Equation~(\ref{eq:erk_wave}) supports traveling wave solutions of the form
$u(\mathbf{r}, t) = U(\mathbf{r} - c\mathbf{n} t)$ for wave speed $c$ and propagation
direction $\mathbf{n}$.

The ZO-1 redistribution dynamics can be coupled to this wave by treating the local
fraction of ZO-1 at apical junctions as a slow structural variable $\phi(\mathbf{r}, t)$
satisfying

\begin{equation}
\frac{\partial \phi}{\partial t} = -k_+(u)\phi + k_-(\phi)(1-\phi),
\label{eq:zo1}
\end{equation}

where $k_+(u)$ is the wave-dependent redistribution rate (increasing in $u$) and $k_-$
is the restoration rate returning ZO-1 to junctions after the wave passes. The coupled
system~(\ref{eq:erk_wave})--(\ref{eq:zo1}) produces the qualitative behavior observed
by Hirano et al.: ZO-1 redistribution following the ERK wave with a delay determined
by the ratio $k_+/k_-$ and the wave speed $c$.

\begin{proposition}
In the coupled system~(\ref{eq:erk_wave})--(\ref{eq:zo1}), the ZO-1 redistribution
front propagates at the same velocity as the ERK wave in the quasi-static limit
$\epsilon \to 0$ of slow recovery dynamics.
\end{proposition}

\begin{proof}
In the quasi-static limit, the recovery variable $v$ varies slowly relative to the
wave transit time, so the ERK wave passes through any local neighborhood $U$ before
$v$ has significantly changed. During the transit, $u$ rises from rest to peak activation
and falls again on the timescale $c^{-1}$ (per unit distance). The ZO-1 variable $\phi$
responds according to~(\ref{eq:zo1}); since $k_+$ is concentrated in the wave region,
$\phi$ decreases during wave passage and recovers at rate $k_-$ afterward. The spatial
locus of minimum $\phi$ therefore propagates with the wave front at velocity $c$.
\end{proof}

\section{ZO-1 Redistribution and Structural Phase Transitions}
\label{sec:phase}

\subsection{Two Organizational Regimes}

The Hirano et al.\ findings identify two distinct organizational regimes for a migrating
epithelial cell: a junction-maintenance regime, in which ZO-1 localizes at apical
tight-junction complexes and the cell participates in maintaining epithelial barrier
integrity and collective cohesion; and a podosome-invasive regime, in which ZO-1
relocalizes to basal podosome-like structures and the cell participates in extracellular
matrix remodeling and directed invasive movement~\cite{Hirano2026}.

These are not merely quantitatively different states along a continuum. They involve
qualitatively different organizational principles: different spatial localization of
key proteins, different cytoskeletal configurations (cortical actin for junction
maintenance versus actin-rich podosome cores), different extracellular interactions
(paracellular barrier maintenance versus matrix degradation and substrate adhesion),
and different contributions to collective behavior (cohesion versus migration).
The transition between them is triggered by the ERK wave and is therefore externally
induced rather than arising from purely internal fluctuation.

\begin{definition}[Organizational Phase]
An organizational phase for a cell in a coupled epithelial monolayer is an equivalence
class of structural configurations with qualitatively similar protein localization patterns,
cytoskeletal architectures, and mechanical interactions with the environment. A phase
transition occurs when the cell moves from one equivalence class to another in response
to a constraint-propagating front.
\end{definition}

The ERK wave is the trigger for the phase transition from junction-main\-te\-nance to
podosome-invasive. The molecular requirements identified by Hirano et al.\ specify the
mechanism: ERK phosphorylation of ZO-1 releases it from junctions (phospho-deficient
mutants fail to redistribute), and the actin-binding and cortactin-binding regions of
ZO-1 are required for podosome localization. The transition is therefore a
phosphorylation-dependent structural reorganization with both energetic and spatial
components.

\subsection{Constraint Density Reallocation}

The RSVP-theoretic interpretation of this transition treats it as a reallocation of
constraint density between organizational modes. In the junction-maintenance phase,
the dominant structural constraint is the maintenance of intercellular cohesion through
tight-junction complexes; ZO-1's contribution to this constraint is local and apical.
In the podosome-invasive phase, the dominant structural constraint shifts to the
generation of directed traction forces and matrix remodeling; ZO-1's contribution
becomes basal and protrusive.

Let $\rho_J(\mathbf{r}, t)$ and $\rho_P(\mathbf{r}, t)$ denote the constraint densities
associated with junction maintenance and podosome activity respectively at position
$\mathbf{r}$ and time $t$. The total constraint density
$\rho_T = \rho_J + \rho_P$ is approximately conserved on the timescale of ZO-1
redistribution, since the total ZO-1 concentration does not change rapidly. The ERK
wave drives a local transfer from $\rho_J$ to $\rho_P$ by enabling the phosphorylation-
dependent redistribution.

\begin{theorem}[Constraint Reallocation]
Under the coupled dynamics~(\ref{eq:erk_wave})--(\ref{eq:zo1}), the passage of an ERK
activation front through a local neighborhood $U$ drives a transient reallocation of
constraint density from junction-maintenance to podosome-invasive modes, with the
magnitude of reallocation proportional to the local ERK peak amplitude $\max_t u(\mathbf{r}, t)$
for $\mathbf{r} \in U$.
\end{theorem}

\begin{proof}[Argument]
The redistribution rate $k_+(u)$ is monotone increasing in $u$ by the phosphorylation
mechanism. Therefore the total ZO-1 displaced from junctions during the wave transit
is
\[
\Delta\phi(U) = \int_0^\infty k_+(u(\mathbf{r}, t)) \phi(\mathbf{r}, t) \, dt,
\]
which is bounded below by $k_+(\max_t u) \cdot \int_0^\infty \phi \, dt$ and above by
$k_+(\max_t u) \cdot \phi(0) \cdot T_{wave}$, where $T_{wave}$ is the wave transit time
through $U$. Both bounds are monotone in $\max_t u$, establishing the claimed proportionality.
The displaced ZO-1 localizes to podosome structures by the actin-binding pathway,
completing the reallocation from $\rho_J$ to $\rho_P$.
\end{proof}

\subsection{Reversibility and Restoration}

The phase transition induced by the ERK wave is not permanent. After the wave passes,
the recovery variable $v$ in equation~(\ref{eq:erk_wave}) suppresses further ERK
activation, and the restoration term $k_-(1-\phi)$ in equation~(\ref{eq:zo1}) drives
ZO-1 back toward junctions. The cell returns to the junction-maintenance phase on a
timescale determined by $k_-$ and the ERK recovery dynamics.

This reversibility is crucial for collective migration. If the phase transition were
irreversible, cells would permanently lose junction-maintenance capacity and the
epithelial sheet would disintegrate. Instead, the wave induces a transient invasive
window: a period during which the cell can participate in directed migration and matrix
remodeling, after which it restores its junction-maintenance function and contributes
to the cohesion that keeps the migrating sheet intact.

The collective consequence is a tissue that can simultaneously maintain structural
integrity and achieve directed migration: different spatial regions are in different
phases at any given time, determined by the local passage of the ERK wave front.
The wave coordinates migration not by bringing all cells into the invasive phase at once
but by propagating the invasive phase through the tissue in a controlled temporal sequence.

\section{Admissibility as a General Principle of Constraint Selection}
\label{sec:admissibility}

The concept of admissibility appears across several independent mathematical domains as
a mechanism for restricting formally possible trajectories, morphisms, or solutions to
those that preserve coherent structural evolution under transformation, perturbation,
or propagation. This recurrence is not terminological coincidence. Before introducing
the RSVP field equations, it is worth surveying these appearances systematically, because
they provide independent mathematical legitimacy for treating accessibility as a
dynamical field variable rather than as a derived property of state evolution. The
novelty of the RSVP framework does not lie in the concept of admissibility itself but
in treating the volume of admissible trajectories as a primary distributed dynamical
quantity rather than as a fixed background constraint.

\subsection{Admissibility in PDE Theory and Weak Solutions}

In the theory of nonlinear hyperbolic partial differential equations, classical solutions
can fail to exist globally due to shock formation and discontinuities. The resulting
weak solutions are generally non-unique; additional selection criteria, called admissibility
conditions, are required to identify physically or organizationally meaningful solutions
among the many that satisfy the equations in the distributional sense. Classical
admissibility conditions include the Lax entropy condition, the viscosity condition,
and the kinetic relation, each embodying a different physical principle for selecting
coherent evolution from formally valid but structurally unstable alternatives.

Recent work has extended admissibility analysis to action rate criteria, formalizing the
selection of weak solutions through the rate at which the system's action functional is
consumed at discontinuities~\cite{Gimperlein2026}. This approach is especially relevant
to the present framework because it treats admissibility as a local-in-time property,
selecting dynamically coherent evolutions at each instant rather than requiring global
minimization. The parallel to RSVP-style constraint-propagating fronts is direct: the
front itself functions as a traveling admissibility criterion, selecting at each spatial
location the organizational regime that preserves coherent structural evolution. The
analysis of hyperbolic equations and non-classical wave structures in related systems
further shows that admissibility conditions govern the emergence of composite waves and
phase boundaries of the type observed in the ERK-ZO-1 transition~\cite{ElAbbassi2025}.

\subsection{Admissibility and Nonuniform Hyperbolicity}

In the theory of dynamical systems, admissibility conditions arise naturally in the
characterization of nonuniform hyperbolicity. A dynamical system exhibits nonuniform
hyperbolicity when the rates of expansion and contraction along trajectories vary in
a measurable but non-constant way, and admissibility conditions for associated operators
are precisely what determine whether the observed trajectory structure is consistent
with a given Lyapunov-exponent profile~\cite{Chemnitz2025}. The key structural feature
is that admissibility in this context is not a property of individual states but of
\emph{weighted trajectory spaces}: the question is not whether a given orbit is possible
but whether the entire family of orbits is organized in a way consistent with the
underlying hyperbolicity structure.

This is directly analogous to the RSVP interpretation of GAD developmental trajectories.
The disorder does not render individual neural configurations impossible; it alters the
organizational geometry of the trajectory family, widening the accessible trajectory
space in ways that are visible at the population level rather than the individual level.
Admissibility in the dynamical-systems sense operates on the same object --- the space
of trajectories rather than the space of states --- as the RSVP accessibility field $S$.

The analysis of admissibility through delay-equation induction offers a complementary
perspective~\cite{Barreira2024}. By embedding lower-dimensional trajectory dynamics
into higher-dimensional delay systems, admissibility conditions acquire a natural
geometric interpretation in terms of the lifted trajectory space, where organizational
constraints become visible as curvature or confinement properties that are not apparent
at the original dimensionality. This is structurally consistent with the RSVP claim
that distributed organizational constraints operating on the trajectory manifold cannot
be reduced to point-wise state properties.

\subsection{Inferential Admissibility and Organizational Coherence}

Admissibility conditions also appear in philosophical logic and inferential semantics
as constraints on which inference rules preserve the coherence of a logical or semantic
system. The classic example is the tonk connective introduced by Prior: a connective
defined by an introduction rule and an elimination rule that, taken together, allow the
derivation of any proposition from any other, collapsing the inferential structure of
the system entirely. Admissibility constraints are what prevent such inferential collapse
by requiring that transition rules preserve the organizational integrity of the
consequence relation~\cite{Topey2026}.

The structural parallel to biological admissibility is not superficial. An inferential
system that loses its admissibility constraints undergoes unbounded expansion of
derivable conclusions --- an explosion of accessible states --- exactly analogous to a
biological system that loses its organizational constraints on trajectory-space accessibility.
In both cases, the failure mode is not incorrect evolution but unconstrained evolution:
the system reaches states that are formally reachable but organizationally incoherent.
The RSVP accessibility field $S$ functions in the biological domain as an admissibility
measure functions in the inferential domain: it bounds the volume of accessible evolution
and thereby preserves the structural coherence that makes organized behavior possible.
This parallel suggests that admissibility is not merely a technical device in each domain
but a general invariant of systems that must maintain coherent organizational structure
under ongoing state evolution.

\subsection{Categorical Admissibility and Structural Preservation}

In category theory and its applications to algebraic topology and localization theory,
admissibility conditions govern which morphisms or localizations preserve relevant
structural properties under functorial transformation. The admissibility of localizations
of crossed modules, for example, determines which localization maps preserve the
higher-dimensional algebraic structure of crossed-module diagrams under
pullback~\cite{Monjon2023}. Admissibility in this setting is a condition on the
transformations themselves rather than on the objects: it distinguishes structure-preserving
mappings from formally valid but structurally destructive ones.

This categorical sense of admissibility connects to the RSVP framework through the
question of which constraint-transport operations preserve organizational coherence
across spatial or temporal projection. The vector field $\mathbf{v}$ in the RSVP
equations does not merely carry activation; it carries the organizational regime through
the tissue. Admissibility, in the categorical sense, would require that this transport
preserve the relevant structural relationships among the $(\Phi, \mathbf{v}, S)$
fields rather than merely relocating activation energy. The multivalued fixed-point
and coincidence-point results that arise in metric fixed-point theory under admissibility
conditions provide further formal grounding for the idea that admissibility restricts
which dynamical correspondences are structurally stable~\cite{Choudhury2025}.

\subsection{Admissibility as a Cross-Domain Invariant}

Across these otherwise disconnected mathematical domains, admissibility repeatedly
emerges as a mechanism for restricting formally possible trajectories to those
preserving coherent structural evolution under transformation, perturbation, or
propagation. Whether the relevant objects are weak PDE solutions, inferential rule
systems, hyperbolic cocycles, localization functors, or collective biological states,
admissibility functions as a constraint-selection principle operating over spaces of
possible organizational evolution. The RSVP framework interprets this recurrence not
as terminological coincidence but as evidence that admissibility is a general invariant
of distributed dynamical organization. The framework's specific contribution is to
treat the volume of admissible trajectories not as a fixed background constraint but
as the primary dynamical variable $S$, a distributed field that evolves under the same
coupled equations as coherence and transport. In the two biological systems examined
here, this dynamical admissibility field is what distinguishes the RSVP interpretation
from both standard diffusion theory and localist entity-centered accounts: the question
is not which states a system can occupy but which \emph{trajectory structures} it can
sustain, and how the volume of those structures evolves under propagating organizational
perturbations.

\section{The RSVP Framework and Constraint Propagation Across Biological Scales}
\label{sec:rsvp}

\subsection{The Field Triple}

The term \emph{field} is used throughout this section in the phenomenological
dynamical-systems sense: a distributed quantity defined over a spatial or
configuration-space domain whose local evolution depends on neighboring states and
transport structure, rather than in the sense of a fundamental gauge or particle field
in high-energy physics. This distinction matters because the RSVP framework is equally
applicable to literal physical space (as in the epithelial monolayer) and to abstract
state manifolds (as in the developmental configuration space relevant to the GAD
findings). In each case, the field triple $(\Phi, \mathbf{v}, S)$ describes the
organizational geometry of the system rather than its physical substrate.

The Relativistic Scalar-Vector Plenum framework characterizes organizational states
through a triple $(\Phi, \mathbf{v}, S)$ defined over an admissible trajectory manifold
$\mathcal{M} \times [0, T]$, where $\mathcal{M}$ is the relevant organizational state
space rather than necessarily a Euclidean domain. In the epithelial migration setting,
$\mathcal{M} \approx \mathbb{R}^2$ is the physical tissue surface. In the neural
developmental setting, $\mathcal{M}$ is the high-dimensional configuration manifold of
structural brain states, which may be reconstructed from imaging or transcriptomic data
via nonlinear dimensionality reduction. In general:
\[
(\Phi, \mathbf{v}, S) : \mathcal{M} \times [0,T] \to \mathbb{R}_{\geq 0} \times
T\mathcal{M} \times [0,\infty),
\]
where $T\mathcal{M}$ denotes the tangent bundle of $\mathcal{M}$, accommodating
transport along the manifold's intrinsic geometry rather than through ambient Euclidean
space. This abstraction unifies what would otherwise appear as domain-specific models
into instances of a single organizational-geometric framework.

\begin{definition}[RSVP Field Triple]
The RSVP organizational state at $(x, t) \in \mathcal{M} \times [0,T]$ is the triple
$(\Phi, \mathbf{v}, S)$ where:
\begin{enumerate}
\item $\Phi(x, t)$ measures local organizational density or coherence: the degree
to which the system at $x \in \mathcal{M}$ is organized around a specific structural
configuration.
\item $\mathbf{v}(x, t) \in T_x\mathcal{M}$ specifies the local transport or propagation
direction and magnitude along $\mathcal{M}$: the direction and rate of constraint
redistribution or signaling flow.
\item $S(x, t)$ measures geometric accessibility: the logarithmic local volume of
dynamically reachable configurations under finite-energy perturbation from the current
state, approximated operationally as
\[
S(x,t) \approx -\log \rho(x,t),
\]
where $\rho(x,t)$ is the local trajectory density on $\mathcal{M}$ at time $t$.
High $S$ implies a large accessible future trajectory bundle and high organizational
plasticity; low $S$ implies a consolidated, narrow trajectory distribution.
\end{enumerate}
\end{definition}

\begin{remark}
The quantity $S$ should not be identified directly with thermodynamic entropy in the
statistical-mechanical sense. It functions as a geometric accessibility measure over
dynamically reachable organizational trajectories. In the biological applications
considered here, $S$ operates at the level of geometric entropy $S_{\mathrm{geom}}$:
the logarithmic volume of admissible trajectory bundles on a reconstructed organizational
manifold. This is distinct from informational entropy $S_{\mathrm{info}}$ (uncertainty
over discrete symbolic states) and from thermodynamic entropy $S_{\mathrm{phys}}$
(microstate counts at the molecular level). Thermodynamic and informational
interpretations are optional specializations rather than foundational commitments of
the framework. The operationalization $S \approx -\log\rho$ is practical precisely
because it makes $S_{\mathrm{geom}}$ inferable from data through inverse local
trajectory density estimation on experimentally reconstructed state manifolds.
\end{remark}

The governing equations couple these three fields over $\mathcal{M}$:
\begin{align}
\frac{\partial \Phi}{\partial t} &= -\mathrm{div}_{\mathcal{M}}(\Phi \mathbf{v})
+ \lambda R(\Phi, S),
\label{eq:phi} \\
\frac{\partial \mathbf{v}}{\partial t} &= -\nabla_{\mathcal{M}}^{\mathbf{v}}\mathbf{v}
- \alpha \,\mathrm{grad}_{\mathcal{M}}\,\Phi
+ \beta \,\mathrm{grad}_{\mathcal{M}}\,S + \mathbf{F}_{ext},
\label{eq:v} \\
\frac{\partial S}{\partial t} &= \kappa \Delta_{\mathcal{M}} S
- \gamma \Phi S + \delta Q(\mathbf{v}),
\label{eq:s}
\end{align}
where $\mathrm{div}_{\mathcal{M}}$, $\mathrm{grad}_{\mathcal{M}}$, and
$\Delta_{\mathcal{M}}$ are the divergence, gradient, and Laplace--Beltrami operators
on $\mathcal{M}$; $R(\Phi, S)$ is a regeneration term encoding how organized structures
rebuild under favorable accessibility conditions; $\mathbf{F}_{ext}$ represents external
forcing; $Q(\mathbf{v})$ represents the accessibility generated by transport flows; and
$\lambda, \alpha, \beta, \gamma, \delta, \kappa > 0$ are coupling constants whose
values are substrate-dependent.

Equation~(\ref{eq:phi}) states that coherence is transported along $\mathcal{M}$ by the
vector field and locally regenerated at a rate depending on existing coherence and
accessibility. Equation~(\ref{eq:v}) states that transport evolves under coherence
gradients (regions of high coherence attract transport) and accessibility gradients
(transport is directed toward regions of higher organizational openness). The
$\beta\,\mathrm{grad}_{\mathcal{M}}\,S$ term encodes the tendency of propagating
organizational dynamics to preferentially expand into regions of higher reorganizational
permissibility, analogous to how activation fronts in excitable media preferentially
propagate through regions below refractory threshold: since $S \approx -\log\rho$,
the accessibility gradient $\mathrm{grad}_{\mathcal{M}}\,S \approx -\mathrm{grad}_{\mathcal{M}}\log\rho$
points away from dense trajectory basins toward regions where organizational
restructuring remains possible. Equation~(\ref{eq:s}) states that accessibility
diffuses over $\mathcal{M}$, is suppressed by high coherence (organized structures close
off alternatives), and is generated by active transport.

A derived quantity of particular interest is the \emph{constraint-flux law}, which
describes how the organizational domain evolves under coupled transport and
accessibility:
\begin{equation}
\partial_t x = F(x, \mathbf{v}, \mathrm{grad}_{\mathcal{M}}\,S), \quad x \in \mathcal{M},
\label{eq:flux}
\end{equation}
capturing the key insight that accessibility gradients shape the directionality of
organizational evolution on $\mathcal{M}$. Because $\mathrm{grad}_{\mathcal{M}}\,S
\approx -\mathrm{grad}_{\mathcal{M}}\log\rho$, this law describes how constraint fronts
preferentially drive organizational evolution away from trajectory concentration basins.
An empirically tractable invariant associated with this law is the ratio
\begin{equation}
\mathcal{I} = \frac{|\mathbf{v}|}{|\partial_t S|},
\label{eq:invariant}
\end{equation}
which measures organizational conversion efficiency: the transport expenditure required
to collapse a unit volume of accessible trajectory space. In the epithelial setting,
$|\mathbf{v}|$ corresponds to wave-front propagation speed and $|\partial_t S|$ to the
rate of junctional ZO-1 departure. In the neural developmental setting, $|\mathbf{v}|$
corresponds to developmental drift velocity on the neural configuration manifold and
$|\partial_t S|$ to attractor-collapse or pruning rate. The hypothesis that
$\mathcal{I}$ takes approximately substrate-invariant values across organizational
systems is a candidate universality claim that would require longitudinal, multi-system
empirical data to evaluate.

Two candidate functional forms for the unspecified terms deserve mention as the most
natural choices consistent with the framework's organizational logic. For the
regeneration term $R(\Phi, S)$, a logistic-exponential form
\begin{equation}
R(\Phi, S) = \Phi(1 - \Phi)\,e^{-S}
\label{eq:R}
\end{equation}
captures the core phenomenology: regeneration is maximal at intermediate coherence
(logistic factor $\Phi(1-\Phi)$, vanishing both at full disorganization $\Phi = 0$
and at full saturation $\Phi = 1$) and suppressed exponentially by high accessibility
(exponential factor $e^{-S}$, since high $S$ corresponds to a plastic, uncommitted
state in which coherent structure cannot easily rebuild). For the accessibility
generation term $Q(\mathbf{v})$, a local kinetic energy density
\begin{equation}
Q(\mathbf{v}) = \tfrac{1}{2}|\mathbf{v}|^2
\label{eq:Q}
\end{equation}
is the most parsimonious choice: transport flows generate accessible future states
at a rate proportional to their squared magnitude, which is consistent with the
Constraint-Flux Law~(\ref{eq:flux}) and with the operationalization
$S \approx -\log\rho$, since fast transport redistributes trajectory density and
thereby raises $S$ in the downstream region. These forms are proposed as minimal
candidates for future parameterization rather than as derivable consequences of the
biological mechanisms.

It should be emphasized that equations~(\ref{eq:phi})--(\ref{eq:flux}) are
phenomenological field equations over an organizational manifold rather than mechanistic
biochemical models. They are intended to capture the organizational logic of constraint
redistribution at a level of abstraction above specific molecular pathways. The equations
are offered as structural scaffolding for interpreting the qualitative dynamics and
generating experimentally discriminable predictions, not as a predictive mechanistic
model awaiting parameter fitting.

\subsection{RSVP Interpretation of the GAD Findings}

In the neural trajectory setting, $\mathcal{M}$ is the high-dimensional manifold of
structural brain configurations; $\Phi(x, t)$ represents the local coherence of
structural neural organization at configuration $x \in \mathcal{M}$ and developmental
time $t$; $\mathbf{v}(x, t)$ represents the developmental drift velocity along
$\mathcal{M}$; and $S(x, t) \approx -\log\rho(x,t)$ represents the geometric
accessibility of the local developmental trajectory bundle, with $\rho$ the local
density of developmentally observed neural configurations.

Healthy development corresponds to a smooth decreasing profile of $S$ along the typical
developmental trajectory: high plasticity in early development, progressive consolidation
in adolescence, stable low plasticity in mature adulthood. The PAD measurement captures
where on this trajectory a given individual's structural configuration falls at their
chronological age.

The GAD population's increased PAD variance, in RSVP terms, corresponds to increased
variance in $S$ across individuals at any given chronological age. The disorder amplifies
inter-individual heterogeneity in the accessibility profile: some individuals consolidate
more rapidly (low $S$, high $\Phi$, positive PAD), others more slowly or atypically
(high $S$, lower $\Phi$, negative PAD). The coupling constant $\gamma$ in
equation~(\ref{eq:s}), which governs the rate at which high coherence suppresses
accessibility, would show elevated inter-individual variance in the GAD population.

\begin{proposition}[RSVP Variance Amplification]
If GAD increases the variance of the coupling parameter $\gamma$ across individuals
without strongly shifting its mean, then the PAD distribution of the GAD population
will satisfy $\mathrm{Var}_{GAD}[\delta] > \mathrm{Var}_{HC}[\delta]$ with
$|\mathbb{E}_{GAD}[\delta] - \mathbb{E}_{HC}[\delta]| \approx 0$, consistent with
the Richier et al.\ findings.
\end{proposition}

The exploratory finding that symptom severity correlates with PAD magnitude after
controlling for medication is consistent with this interpretation: symptom severity
reflects the magnitude of the perturbation to $\gamma$, which drives the individual
toward either extreme of the accessibility distribution. High symptom severity
corresponds to large $|\xi_i|$, pushing the individual toward either early consolidation
or prolonged plasticity.

\subsection{RSVP Interpretation of the ERK Wave Dynamics}

In the tissue migration setting, $\mathcal{M} \approx \mathbb{R}^2$ is the physical
epithelial surface; $\Phi(x, t)$ represents local organizational coherence of the
junction-maintenance architecture (tight junctions, apical ZO-1 localization, cortical
actin organization); $\mathbf{v}(x, t)$ represents the local transport of signaling
activation (ERK wave propagation direction and velocity along the tissue surface); and
$S(x, t) \approx -\log\rho(x,t)$ represents local structural accessibility, with $\rho$
the local density of observed cellular configurations at $(x,t)$, with high $S$ enabling
the transition to podosome-invasive organization and low $S$ enforcing junction-maintenance.

The resting epithelial monolayer corresponds to high $\Phi$ (well-organized junctions),
low $\mathbf{v}$ (no directed transport), and low $S$ (junction maintenance is the dominant
organizational attractor, podosome formation is suppressed). The ERK wave corresponds to
a propagating region of elevated $\mathbf{v}$, which through equation~(\ref{eq:s})
generates elevated $S$ in its wake (transport flows open up accessibility). This elevated
$S$, through the term $-\gamma \Phi S$ in equation~(\ref{eq:s}), reduces $\Phi$ locally,
enabling the regime transition to podosome-invasive organization.

After the wave passes, $\mathbf{v}$ returns to low amplitude, $S$ decays through the
$\kappa \nabla^2 S$ diffusion term (accessibility spreads and diminishes), and the
regeneration term $\lambda R(\Phi, S)$ in equation~(\ref{eq:phi}) drives $\Phi$ back
toward its junction-maintenance value. The tissue restores junction-maintenance
organization while the wave front continues propagating through downstream regions.

The RSVP framework thus captures the key qualitative features of the Hirano et al.\
system: wave-triggered regime transitions, ZO-1 redistribution as a transient suppression
of $\Phi$, and restoration to junction-maintenance after wave passage.

\subsection{RSVP Interpretation of the GBF Neuroimaging Framework}

The Wang et al.\ geometric basis function framework provides what is in some respects
the most direct empirical instantiation of the RSVP manifold formalism~\cite{Wang2026}.
In that paper, $\mathcal{M}$ is not an abstract organizational manifold but a
literal Riemannian surface: the individual's cortical mesh reconstructed from structural
MRI and endowed with the intrinsic metric of the cortical geometry. The Laplace--Beltrami
operator $\Delta_{\mathcal{M}}$ is computed explicitly on this surface, yielding
eigenmodes $\psi_i$ that are precisely the harmonic basis of the manifold. Neural
activity is then expressed as $x(t) = \sum_i \theta_i(t) \psi_i$, a decomposition
that is RSVP-compatible at the level of notation: the coherence field $\Phi$
corresponds to the spatially organized amplitude distribution $x(t)$ over $\mathcal{M}$,
and the eigenmode coefficients $\theta_i(t)$ carry the temporal dynamics.

The spectral prior $\Sigma^{-1} = \mathrm{diag}(-\beta/\log\lambda_i)$ is the most
directly interpretable component from the RSVP accessibility perspective. Because
$\lambda_i$ increases with spatial frequency, $-\beta/\log\lambda_i$ decreases with
$\lambda_i$, assigning smaller prior variance (higher regularization) to
high-frequency modes. Equivalently, high-frequency eigenmodes are penalized not
because they are physically forbidden but because they correspond to high-curvature,
short-wavelength configurations that are dynamically inaccessible under realistic
signal-to-noise and physiological constraints. In RSVP terms, this prior is an
admissibility weighting: $S(\psi_i) \propto -\log\lambda_i$, assigning lower
accessibility to configurations that oscillate rapidly over the cortical surface.
The logarithmic form is significant: it matches the RSVP operationalization
$S \approx -\log\rho$ if the local eigenmode density on the spectral manifold is
taken to scale as $\rho \propto \lambda_i$ in the high-frequency regime, consistent
with the Weyl law for Laplace--Beltrami eigenvalue asymptotics.

The finding that approximately 200--300 modes suffice to reconstruct whole-brain
dynamics with high fidelity is directly interpretable as an empirical upper bound on
the dimensionality of the admissible trajectory manifold $\mathcal{M}_{adm}$. Of the
approximately 32,000 spatial degrees of freedom in the cortical vertex representation,
only a low-dimensional submanifold is dynamically admissible under the cortical
geometry prior. The remaining dimensions are not activated because they are
admissibility-suppressed: their high curvature places them in regions of low $S$,
making them dynamically inaccessible rather than merely unexplored. This is
structurally identical to the RSVP prediction that organized biological systems
maintain low effective dimensionality by suppressing high-$S$ organizational modes
through the coupling term $-\gamma \Phi S$ in equation~(\ref{eq:s}).

The optical-flow wave propagation analysis in the Wang et al.\ paper maps directly
onto the RSVP vector field $\mathbf{v}$. Phase gradients across the cortical surface
define a local velocity field
\[
\mathbf{v}(\mathbf{r}, t) = -\frac{|\partial_t \phi|}{|\nabla \phi|^2}\,\nabla\phi,
\]
which is precisely a surface-intrinsic transport field of the type $\mathbf{v} \in
T\mathcal{M}$ in the RSVP formalism. The convergence hubs identified by DBSCAN
clustering of streamline endpoints correspond to attractor basins in the $\Phi$ field:
regions of high coherence density where transport trajectories terminate. This gives
the invariant $\mathcal{I} = |\mathbf{v}|/|\partial_t S|$ a direct operationalization
in the GBF framework: $|\mathbf{v}|$ is the optical-flow propagation speed (reported
in mm/ms), and $|\partial_t S|$ can be approximated by the rate of change of the
effective mode dimensionality as the wave passes through a cortical region.

The GBF paper therefore functions simultaneously as an empirical confirmation of the
RSVP manifold geometry and as a methodological blueprint for operationalizing the
abstract fields $(\Phi, \mathbf{v}, S)$ in neural data. The cortical manifold is the
organizational substrate; the eigenmodes are the admissible trajectory basis; the
spectral prior is the accessibility field; and the propagating cortical waves are
constraint-propagating fronts carrying organizational transitions across the cortical
surface. The representational framing of the GBF paper --- cortical geometry as a
prior constraining inverse solutions --- is, from the RSVP perspective, a special
case of a more general claim: cortical geometry is constitutive of the admissibility
structure of neural trajectories, not merely a useful regularization device.

\subsection{Cross-Scale Structural Invariants}

The RSVP interpretation of all three systems reveals a shared set of structural invariants
that persist across substantial differences in biological scale, mechanism, and measurement
modality.

The first invariant is the primacy of constraint propagation over signal diffusion.
In all three systems, the causally relevant dynamics involves the propagation of
regime-enabling fronts or admissibility-selecting operators rather than the diffusion
of quantitative signals. ERK waves enable qualitative phase transitions in the tissue;
developmental perturbations in GAD alter the qualitative topology of accessible
trajectory space; cortical eigenmodes select admissible neural configurations by
suppressing high-curvature trajectories through the spectral prior. In no case is the
dominant organizational signal a concentration gradient or a mean-field intensity;
it is either a traveling boundary between dynamical regimes or an admissibility
operator that restricts the accessible configuration space.

The second invariant is variance amplification as the primary signature of altered
organizational dynamics, rather than mean-field displacement. Increased PAD variance in
GAD and increased heterogeneity in wave propagation velocity and ZO-1 redistribution
magnitude across individual cells are both expressions of this pattern. The GBF paper
provides the complementary case: reconstruction quality saturates at approximately
300 modes and degrades when higher-order modes are included, because expanding the
active dimensionality beyond the admissible manifold amplifies ill-conditioning and
noise rather than adding signal. The systems do not shift their average behavior;
they are governed by the geometry of the admissible distribution, and it is departures
from this geometry — whether through pathological dilation in GAD, constrained
ERK wave propagation in migration, or inadmissible mode inclusion in source imaging ---
that constitute the observable signatures of organizational disruption or miscalibration.

The third invariant is the identification of reversibility with organizational competence.
In all three systems, the capacity to undergo regime transitions and return to baseline
functions as a marker of health. GAD involves a loss of developmental trajectory focus
without a corresponding loss of the mean trajectory; disrupted ERK-ZO-1 coupling,
whether through MEK inhibition or phospho-deficient ZO-1 mutants, impairs not the static
structure of junctions but the dynamic capacity for coordinated migration. In the GBF
setting, the logarithmic spectral prior achieves precisely this: it permits moderate
excursions into higher spatial frequencies (reversible activation of finer-grained modes)
without permanently committing to configurations outside the admissible manifold.
The prior moderates penalization rather than truncating modes absolutely, preserving
the capacity for controlled excursion and return.

The fourth invariant is distributed coordination without centralized control. In none
of the three systems does a central controller direct the macroscopic behavior. Neural
developmental trajectories are shaped by distributed gene expression, environmental
input, and accumulated experience; collective tissue migration is shaped by distributed
ERK wave propagation and intercellular mechanical coupling; cortical source dynamics
are shaped by the distributed eigenmode structure of the cortical geometry itself,
with no single region or module imposing the admissibility constraint. Coherent global
behavior emerges from local dynamics and their spatial coupling through the manifold
geometry, not from any supralocal supervisory process.

\section{Implications for Psychiatry, Development, and Collective Intelligence}
\label{sec:implications}

\subsection{Psychiatric Taxonomy and Field Geometry}

The GAD brain-age findings suggest a reconsideration of psychiatric taxonomy that has
broader implications. Standard diagnostic categories in psychiatry implicitly assume
that conditions correspond to relatively discrete regions of neurobiological state space,
distinguishable from health and from each other by characteristic mean-field features.
The accumulating evidence from large-scale neuroimaging consortia increasingly challenges
this assumption.

The RSVP framework suggests an alternative: psychiatric conditions correspond primarily
to deformations of the accessibility geometry of developmental trajectory space rather
than to specific locations in neural configuration space. A condition may manifest
predominantly as trajectory-variance amplification (GAD, perhaps; also consistent with
the enormous heterogeneity reported across autism spectrum studies), or as trajectory-
compression into a narrow pathological attractor (some features of OCD or addiction),
or as loss of trajectory-regeneration capacity (aspects of treatment-resistant depression).

This reframing has direct implications for biomarker development. If the primary
neurobiological signal is a geometric deformation of trajectory space rather than a
mean-field configuration shift, then effective biomarkers must capture distributional
properties of neural trajectories rather than point estimates of brain state. The CNN
approach used by Richier et al.\ represents a step in this direction, but even more
powerful would be methods that directly estimate the local dimensionality, curvature,
or accessibility volume of the individual's trajectory space from longitudinal imaging data.

\subsection{Collective Intelligence and Wave-Mediated Coordination}

The ZO-1 and ERK findings have implications beyond wound healing and embryonic development.
They contribute to an emerging picture of tissues as genuinely computational systems
that solve coordination problems through distributed wave-mediated signal processing
rather than through cellular autonomy or central control.

This picture is continuous with work on morphogenetic fields~\cite{Levin2012},
bioelectric coordination~\cite{Maturana1980}, and the general theory of self-organization
in biological systems~\cite{Camazine2003}. The RSVP framework positions these
observations within a common mathematical language that facilitates comparison across
systems and scales. The propagating constraint front is a mathematical object that
appears both in neural developmental dynamics and in epithelial collective migration;
recognizing this shared structure opens possibilities for cross-domain theoretical
transfer.

In particular, the wave-mediated coordination mechanism raises the question of whether
similar constraint-propagating fronts play organizational roles in systems traditionally
modeled without reference to wave dynamics: neural circuits during learning, immune
system coordination, and social coordination in human groups. The RSVP framework
predicts that wherever distributed coordination is achieved without central control,
constraint-propagating fronts of some form will be identifiable. It should be
emphasized that this essay does not claim these systems share mechanistic identity
with the ERK wave or neural developmental systems discussed here; only that they may
instantiate homologous organizational invariants, in the sense that their dynamical
organization may be governed by structurally similar admissibility constraints operating
over analogous trajectory spaces. Whether this homology is deep or shallow is an
empirical question that requires system-specific analysis.

\subsection{Developmental Timing and Phase Transition Windows}

Both papers implicate temporal dynamics as a crucial organizational variable. The GAD
findings show that the increased PAD variance is more pronounced in adults over
twenty-five than in younger individuals, suggesting that the developmental timing of
GAD onset and consolidation shapes the neurological trajectory in ways that diverge
over time. The ERK wave findings show that the phase transition from junction-maintenance
to podosome-invasive organization occurs in a temporally precise window defined by
the wave transit time.

The concept of a constraint-propagating front provides a unified language for both:
developmental trajectories pass through constraint-propagating fronts associated with
maturational transitions (critical periods, hormonal reorganization, experience-dependent
consolidation), and the organizational consequences of these fronts depend on the
system's admissibility structure at the time of passage. GAD may alter the admissibility
structure in ways that amplify the divergence of trajectories at subsequent constraint-
propagating fronts, producing the increasing variance observed in adults.

\subsection{Extracellular Matrix Remodeling as Environmental Constraint Writing}

A further implication of the Hirano et al.\ findings concerns the bidirectional nature
of the interaction between cell and environment. Podosomes are not merely sensors of
the extracellular matrix; they actively remodel it through proteolytic degradation and
physical deformation. The migrating tissue therefore continuously rewrites the substrate
it inhabits while simultaneously reorganizing itself in response to that substrate.

This bidirectional coupling is a specific instance of the general principle that
organized systems do not merely respond to environmental constraints but modify them
through their activity. In RSVP terms, the vector field $\mathbf{v}$ does not merely
propagate through a fixed domain $\Omega$ but through a domain whose admissibility
structure is itself being modified by the system's activity. This self-referential
modification of the constraint environment is a general feature of living systems
and represents a fundamental limitation of any modeling framework that treats the
environment as a fixed external input.

\subsection{Relation to Existing Theoretical Frameworks}

The RSVP interpretation developed here does not emerge from a theoretical vacuum, and
situating it relative to adjacent frameworks clarifies both what it contributes and
where its claims remain genuinely distinctive.

The closest established framework is Friston's Free Energy Principle~\cite{Friston2010},
which also treats biological systems as minimizing a scalar quantity (variational free
energy) over generative model parameters and sensory states. The RSVP framework shares
with the Free Energy Principle the emphasis on distributed inference and the rejection
of passive stimulus-response causation, but differs in its explicit treatment of
spatial field dynamics: the Free Energy Principle does not naturally accommodate
propagating constraint fronts or spatially distributed phase transitions without
additional architectural assumptions. The ERK wave system, where the organizational
transition is intrinsically spatial and wave-mediated, is more naturally described
by a field-theoretic formalism than by a variational inference framework applied locally.

Morphogenetic field theory in the tradition of Turing~\cite{Turing1952} and its
contemporary descendants in bioelectric field research~\cite{Levin2012} also address
spatially distributed coordination in biological systems. Reaction-diffusion systems of
the Turing type produce spatial pattern through local activation and long-range inhibition,
and the ERK wave system has features consistent with this class of dynamics. The
distinction is that standard Turing pattern formation produces static spatial patterns
as attractors, whereas the ERK system produces traveling waves that transiently reorganize
structure rather than selecting a fixed spatial configuration. The constraint-propagating
front formalism is intended to capture this transient, trajectory-level reorganization
that static pattern-formation theory does not naturally address.

Active matter physics provides a related set of tools for understanding collective
dynamics in epithelial tissues, treating cell monolayers as oriented fluids with
internal driving forces~\cite{Haken1983}. Active matter descriptions successfully account
for many mechanical features of collective migration but typically do not incorporate
the signaling dynamics that drive phase transitions between organizational regimes.
The RSVP framework treats the signaling and mechanical layers as coupled through the
$(\Phi, \mathbf{v}, S)$ triple, which is complementary to rather than competitive with
active matter descriptions.

Dynamical systems approaches to psychiatry, associated with work in the tradition of
Kelso~\cite{Kelso1995} and more recent computational psychiatry, also treat mental
states as attractors in high-dimensional state spaces and pathology as altered attractor
geometry. The RSVP framework's treatment of GAD as a widening of the admissibility
distribution is continuous with this tradition but adds the specific claim that the
dominant signature is variance amplification in trajectory space rather than attractor
displacement, a distinction that has direct empirical consequences for biomarker design.

The RSVP approach is perhaps most distinctive in its insistence on a unified field-
theoretic language that applies across biological scales without requiring scale-specific
mechanisms. Whether this cross-scale applicability is a theoretical virtue or a liability
(by sacrificing mechanistic specificity) is a legitimate methodological question that
future empirical work will need to address.

\section{Conclusion: Biology as Dynamic Admissibility Geometry}
\label{sec:conclusion}

The two papers considered in this essay resist explanation by standard entity-centered
biological models. In the case of GAD, the relevant neurobiological signal is not the
mean configuration of the brain but the variance of its developmental trajectories;
in the case of collective epithelial migration, the relevant organizational phenomenon
is not the behavior of individual cells but the propagation of constraint-redistributing
waves across a coupled tissue medium. Both findings reveal that the appropriate
descriptive level for understanding the macroscopic behavior of complex biological
systems is the geometry of their organizational dynamics: the structure of accessible
trajectory spaces, the propagation of constraint-enabling fronts, and the phase
transitions between organizational regimes.

The RSVP framework, as a generalized field theory of organizational dynamics, provides
a common mathematical language for interpreting both systems. The scalar coherence
field $\Phi$, the vector transport field $\mathbf{v}$, and the entropic accessibility
measure $S$ together characterize the organizational state in a way that captures
the distributional and dynamical properties that scalar biomarker approaches lose.
The coupled field equations reproduce the qualitative behavior of both systems and
generate specific structural predictions: that variance amplification is a primary
signature of trajectory-space deformation in GAD; that ERK waves drive constraint
reallocation from junction-maintenance to podosome-invasive organization; that
organizational health corresponds to reversible phase transition capacity rather
than static structural configuration.

The deeper implication is methodological. These findings expose a limitation in
strongly localist approaches to biological explanation: entity-centered frameworks,
built around the objects identified by structural biochemistry, molecular genetics, and
histology, are productive for many problems but become inadequate precisely where
distributed temporal dynamics and collective coordination are primary. The brain-age
paper's discovery that diagnosis fails to predict individual trajectory and that deep
multivariate networks outperform univariate biomarkers is not an anomaly; it is a signal
that the organizational structure of interest is carried at a level that scalar summaries
cannot reach. The ZO-1 paper's discovery that a structural protein participates in
distributed wave-mediated coordination rather than serving merely as a static barrier
component is not an exception; it is an instance of a general principle.

That principle is the one identified here: biological systems at multiple scales achieve
organizational coherence through the propagation of constraint-redistributing dynamics
across coupled fields, not through the aggregation of fixed object properties. Formalizing
this principle with mathematical precision, as the RSVP framework attempts, may become
increasingly important for the development of theoretical biology as the explanatory
targets shift from component identification toward distributed organizational dynamics.

\section{Conceptual Clarifications and Theoretical Boundaries}
\label{sec:clarifications}

The organizational geometry developed in preceding sections is ambitious enough in
scope that several conceptual risks require direct address. The following subsections
clarify what the framework does and does not claim, and sharpen the distinctions that
are most likely to be blurred in interdisciplinary reading.

\subsection{Constraint versus Cause}

The RSVP field triple $(\Phi, \mathbf{v}, S)$ could be misread as proposing hidden
causal substances or vitalist organizing forces that drive biological systems from
outside their physical substrate. This reading would be incorrect. The framework
proposes nothing of the kind.

The RSVP fields are organizational constraint structures: mathematical representations
of which organizational transitions are accessible, at what rates, and in which
directions, given the system's current state and coupling geometry. They are not
energetic fluids, hidden forces, or ontologically independent entities. They are
closer in spirit to the reachability sets of control theory, the feasibility regions
of constrained optimization, or the admissible trajectory bundles of dynamical
systems analysis --- all of which are mathematical objects describing the geometry
of possible evolution rather than forces producing that evolution.

The distinction matters practically. When the vector field $\mathbf{v}$ propagates
through the epithelial tissue, the RSVP interpretation does not claim that $\mathbf{v}$
is causing the ERK wave; it claims that $\mathbf{v}$ is the field-theoretic
representation of the organizational consequences of the wave at the level of
constraint geometry. The ERK wave has entirely ordinary molecular causes: receptor
activation, MAPK cascade propagation, paracrine signaling. What the RSVP framework
describes is how these mechanistic events reshape the landscape of accessible
organizational transitions at the tissue scale. Mechanism and constraint geometry
are complementary descriptions operating at different levels of abstraction, not
competing explanations of the same phenomenon.

\subsection{Why Variance Is Not Noise}

The treatment of PAD variance as the primary organizational signal rather than as
noise around a mean effect requires explicit justification, because the biomedical
convention is precisely the opposite: variance is typically treated as a nuisance
to be controlled or averaged away, leaving the mean effect as the quantity of
theoretical interest.

The RSVP framework reverses this priority in a principled way. Variance in an
accessibility field $S$ directly reflects the geometry of the admissibility manifold:
high variance corresponds to a population whose members occupy structurally divergent
positions in the organizational trajectory space, while low variance corresponds to
a population concentrated near a common trajectory. When GAD produces increased PAD
variance without a corresponding mean shift, this is not statistical noise; it is
evidence that the disorder has deformed the admissibility geometry without translating
the population mean.

This reversal is well-supported in adjacent literatures. At critical phase transitions
in physical systems, variance diverges before the transition while the mean may remain
stationary; the variance signal is the leading indicator of organizational change.
In heteroscedastic dynamical systems, increased variance can signal proximity to a
bifurcation boundary where small perturbations produce qualitatively different
outcomes depending on initial condition. In stochastic resonance, variance in the
input is not noise but the mechanism by which suprathreshold organization is achieved.

The RSVP framework generalizes this insight: in any system governed by propagating
admissibility constraints, the variance of observable state variables across system
components should be treated as primary organizational data rather than statistical
residual. This has immediate implications for psychiatric biomarker design. If the
primary signal of GAD is distributional rather than directional, then mean-based
biomarkers will systematically fail to detect it, and variance-based measures ---
distributional spread of PAD, trajectory substructure in longitudinal cohorts,
inter-individual heterogeneity in accessibility profiles --- become the appropriate
clinical targets.

\subsection{Operationalizing the Accessibility Field}

The accessibility field $S \approx -\log\rho$ is conceptually grounded but
experimentally challenging to measure. Beyond the basic inverse-density operationalization,
several candidate estimators are available depending on the biological system and data type.

For single-cell transcriptomic data, the local intrinsic dimensionality (LID) of
the cell state distribution provides a natural proxy: high LID regions correspond
to high $S$ (many accessible nearby states), while low LID regions correspond to
committed, low-$S$ attractor basins. Methods from the manifold learning literature
--- including TWO-NN estimators, persistent homology density, and diffusion-map
volume measures --- can operationalize this directly from high-dimensional state
vectors without specifying the manifold embedding in advance.

For imaging data of the type used in the GBF and GAD papers, the local
mode-participation ratio of the spectral decomposition serves as a proxy: regions
where many eigenmodes contribute significantly (high mode participation) correspond
to high $S$, while regions dominated by a single low-frequency mode correspond to
low $S$. The GBF paper's finding that approximately 300 modes suffice for reconstruction
implies that the mode-participation ratio saturates sharply at this threshold, marking
the effective boundary of the admissible manifold.

For collective migration data, the local curvature of the traction-force field and
the local velocity divergence of the cell flow both provide indirect measures of $S$:
high divergence and high curvature correspond to regions where the tissue is undergoing
regime transition, while low divergence corresponds to consolidated migration in a
low-$S$ state. Graph Ricci curvature, applied to the cell contact network, also
provides a substrate-independent proxy for local organizational openness.

For neural systems more broadly, local Lyapunov exponent spectra estimated from
neural time series provide access to the effective dimensionality of the accessible
trajectory bundle, with positive exponents corresponding to high $S$ and convergent
dynamics corresponding to low $S$.

None of these estimators are equivalent to the theoretical $S$, and each introduces
its own assumptions and measurement artifacts. The appropriate choice is
substrate-dependent, and the RSVP framework does not privilege any one. The
common feature is that all estimate the local volume of dynamically reachable
configurations from the observed distribution of trajectories rather than from
first-principle microstate enumeration.

\subsection{Traveling Fronts versus Equilibrium States}

A recurrent theme across all three empirical systems is that the organizationally
dominant structures are not stable equilibria but transient propagating boundaries
between regimes. This is a substantial departure from the equilibrium-centered
paradigm that dominates much of theoretical biology, where the appropriate question
is "what attractor does the system settle into?" rather than "how does the transition
between attractors propagate?"

The RSVP framework treats propagating fronts as primary for the following reason.
A stable attractor is a region of low $S$: the system has committed to a narrow
trajectory bundle and is not easily perturbed out of it. The dynamically interesting
events --- development, learning, disease, collective migration, cortical reorganization
--- are not characterized by life in the attractor but by transitions between
attractors. These transitions propagate through the tissue or the developmental
manifold as moving boundaries between high-$S$ and low-$S$ regions, and it is the
topology, velocity, and reversibility of these boundaries that carry the organizationally
critical information.

This reframing has empirical consequences. Classical reaction-diffusion theory explains
biological pattern formation through the emergence of stable spatial attractors (Turing
patterns, standing waves). The RSVP framework is not competing with this account for
static pattern formation; it is offering a different account for systems whose primary
organizational signature is the traveling front rather than the stable pattern. The ERK
wave does not produce a Turing-type stable spatial organization; it produces a transient
reorganization that propagates and then resolves. The GBF cortical wave is not a static
network pattern; it is a millisecond-scale propagating activation whose trajectory
determines downstream coordination. The PAD trajectory is not a fixed developmental
attractor; it is a dynamic path whose variance structure captures individual differences
in developmental commitment speed.

In each case, the front is the object of theoretical interest, and the equilibrium
is merely the region the system inhabits between fronts. Designing biological theories
and biomarkers around equilibria systematically loses the most organizationally
informative events.

\subsection{Admissibility Compression versus Destructive Reduction}

The GBF finding that 200--300 eigenmodes reconstruct whole-brain dynamics to high
fidelity from 32,000 cortical vertices might be read as saying that the brain is
"simple" or that the high-dimensional data is redundant. This reading misses the
critical distinction between admissibility-preserving compression and destructive
reduction.

Destructive reduction collapses a high-dimensional system to a low-dimensional
representation by discarding the information that distinguishes relevant organizational
configurations from each other. Averaging across a population, binning continuous
variables, or projecting to the first principal component all produce reductions that
lose structural information. The result is a lower-dimensional description that may
be easier to analyze but that no longer supports the organizational distinctions
the original system could make.

Admissibility-preserving compression is structurally different. It identifies the
low-dimensional submanifold $\mathcal{M}_{adm}$ within the full configuration space
$\mathcal{M}$ that is consistent with the system's organizational constraints, and
represents the system's behavior as dynamics on $\mathcal{M}_{adm}$ rather than
on $\mathcal{M}$. The remaining dimensions are not discarded because they are
irrelevant; they are suppressed because the system's admissibility structure prevents
them from being accessed under realistic operating conditions. Including them in
analysis does not add information; it adds noise.

The Wang et al.\ spectral prior $\Sigma^{-1} = \mathrm{diag}(-\beta/\log\lambda_i)$
implements precisely this distinction. It does not truncate modes above a fixed index;
it weights all modes by their geometric admissibility, suppressing high-frequency
configurations smoothly in proportion to their spectral energy. The result is an
admissibility-preserving projection that retains the full organizational structure
of the admissible manifold while numerically suppressing the inadmissible dimensions.
The 300-mode saturation threshold is not evidence that the brain's dynamics are
simple; it is evidence that the brain's admissibility geometry is approximately
300-dimensional within the cortical manifold, a substantial but constrained organizational
capacity.

This distinction also applies to the GAD case. A mean-field biomarker that averages
PAD across the GAD population is destructively reductive: it compresses the population
distribution to its center of mass and loses the variance geometry that carries the
primary organizational signal. An admissibility-preserving analysis would instead
characterize the shape of the trajectory-space distribution --- its variance, its
multimodal substructure, its dependence on developmental timing --- and treat this
shape as the primary data. The RSVP framework is in part an argument for replacing
destructive reduction with admissibility-preserving projection as the default
analytical strategy for complex biological systems.

\subsection{Failure of Reversibility and Organizational Collapse}

The interpretation developed throughout this essay treats reversible phase-transition
capacity as a defining feature of organizational competence. Healthy systems are not
those that remain fixed within a single attractor basin, but those capable of entering
temporary reorganizational regimes and subsequently restoring structural coherence after
the perturbation has passed. This raises the corresponding question: what occurs when
restoration dynamics fail?

The RSVP framework predicts that irreversible organizational collapse occurs when the
rate of coherence regeneration becomes persistently insufficient to counteract the
expansion of accessibility generated by destabilizing transport dynamics. In terms of
the coupled field equations, this corresponds qualitatively to regimes in which
\[
\lambda R(\Phi, S) < \gamma \Phi S
\]
over extended intervals, such that the suppressive effect of elevated accessibility on
coherence exceeds the system's regenerative capacity. Under these conditions, the system
no longer undergoes reversible excursions between organizational phases but becomes
trapped in pathological attractor geometries. The resulting dynamics are not merely
unstable but structurally self-reinforcing: reduced coherence elevates accessibility
variance, elevated accessibility further impairs coherence restoration, and the system
progressively loses the capacity to return to its prior organizational basin.

This general structure appears across multiple biological domains. In epithelial systems,
persistent failure of restoration dynamics may produce irreversible junctional breakdown,
fibrotic remodeling, or metastatic escape, where transient invasive regimes cease to
resolve back into cohesive tissue architecture. In neural systems, analogous dynamics
may underlie forms of pathological attractor locking associated with chronic stress,
treatment-resistant psychiatric states, or neurodegenerative consolidation processes,
where developmental flexibility collapses into rigid low-dimensional trajectories from
which recovery becomes increasingly improbable.

The framework therefore predicts that pathological irreversibility should not primarily
be identified by elevated perturbation magnitude alone, but by the failure of restoration
kinetics relative to accessibility expansion. Two systems exposed to similar perturbation
energy may diverge sharply in long-term outcome depending on whether regenerative
constraint reconstruction remains dynamically viable after perturbation passage.

This interpretation also clarifies why many biological pathologies exhibit hysteresis.
Once the admissibility geometry has been sufficiently deformed, simply removing the
original perturbation may be insufficient to restore the prior organizational state.
The system must actively reconstruct the lost coherence field, often requiring external
intervention, regenerative scaffolding, or prolonged low-perturbation intervals before
the original trajectory manifold becomes accessible again. The field-memory prediction
(Prediction~4 in Section~\ref{sec:predictions}) is the empirically tractable near-term
expression of this deeper claim about irreversibility thresholds and restoration kinetics.

\subsection{Structural Homology versus Metaphorical Analogy}

A persistent risk in cross-domain theoretical work is the substitution of metaphorical
similarity for genuine structural correspondence. Because the RSVP framework employs a
common mathematical language across neural development, epithelial migration, and
geometry-constrained cortical dynamics, it is necessary to clarify the level at which
the claimed unification operates.

The framework does not claim that these systems share mechanistic identity. ERK waves,
developmental neural trajectories, and cortical eigenmode dynamics are implemented
through distinct substrates, governed by different molecular mechanisms, and measured
through different experimental modalities. The RSVP interpretation does not erase
these differences; it operates at a different level of description. The claim instead
concerns structural homology at the level of organizational dynamics: all three systems
exhibit propagating transitions between admissibility regimes rather than static
equilibrium organization, variance amplification as the dominant organizational signature,
distributed coordination without centralized control, reversible phase-transition
competence as a marker of organizational health, and constrained evolution on
low-dimensional admissible manifolds embedded within higher-dimensional configuration
spaces.

These are not merely linguistic similarities, because they generate experimentally
distinguishable predictions. The framework predicts transient widening and reconvergence
of PAD variance during developmental reorganization, topology-sensitive migration
coherence under wave-structured ERK activation, hysteretic modification of redistribution
thresholds after repeated wave exposure, and dimensionality saturation in
admissibility-constrained neural reconstruction. A purely metaphorical analogy would
not constrain empirical expectation in these ways. The RSVP framework therefore stands
or falls not on whether the systems appear similar under broad interpretation, but on
whether the predicted invariants are experimentally recoverable across organizational
domains.

The distinction is methodologically important. Biological science has often relied on
productive metaphors --- the genome as code, the brain as computer, the immune system
as defense network --- but metaphor alone does not establish mathematical continuity.
The present framework instead proposes that admissibility geometry constitutes a
cross-domain organizational invariant whose recurrence reflects shared structural
constraints on distributed dynamical systems rather than rhetorical resemblance.
Whether this claim ultimately reflects deep organizational universality or merely a
high-level abstraction common to many sufficiently complex systems remains an open
empirical question. The framework gains scientific legitimacy precisely to the degree
that its proposed invariants survive increasingly substrate-specific experimental tests.

\subsection{Admissibility Suppression and Computational Tractability}

One of the most important implications of the geometry-aware neuroimaging framework
is that biological systems may remain computationally manageable not because their
underlying state spaces are intrinsically small, but because the overwhelming majority
of formally possible configurations are dynamically inadmissible.

The cortical reconstruction problem examined by Wang et al.\ begins with approximately
32,000 spatial degrees of freedom, yet the empirically admissible dynamics are captured
to high fidelity by only 200--300 geometric basis functions. This does not imply that
the remaining dimensions are irrelevant in principle; rather, the spectral prior
suppresses configurations whose geometric curvature places them outside the dynamically
accessible manifold under physiological operating conditions. The RSVP framework
generalizes this observation: organized systems maintain tractability through
\emph{admissibility suppression}, in which the active trajectory bundle occupies only
a narrow submanifold $\mathcal{M}_{adm} \subset \mathcal{M}$ of the full configuration
space. The accessibility field $S \approx -\log\rho$ functions precisely as a measure
of this suppression geometry, with regions of low $S$ corresponding to concentrated
admissible trajectories and therefore to reduced effective computational dimensionality.

This perspective reframes several otherwise disconnected observations across biology
and cognition. Neural systems exhibit sparse activation despite enormous combinatorial
capacity; developmental systems reliably converge despite high-dimensional molecular
variation; collective tissues coordinate coherently despite local stochasticity. In each
case, admissibility constraints dynamically suppress most formally possible organizational
trajectories, allowing coherent large-scale behavior to emerge without requiring
exhaustive exploration of the underlying state space. The resulting organizational
economy is structurally analogous to manifold learning and compressed representation
in machine learning: high-dimensional observational spaces often contain low-dimensional
dynamical structure embedded within them, and successful coordination depends on
identifying the admissible substructure rather than representing every formal degree
of freedom equally.

This interpretation also clarifies why perturbations near admissibility boundaries are
especially destabilizing. As the system approaches regions where the effective dimension
of the admissible manifold expands rapidly, the computational burden of maintaining
coherent coordination increases sharply. Variance amplification, trajectory divergence,
and loss of reversibility emerge naturally under these conditions because the system is
forced to navigate a substantially larger accessible trajectory bundle. The framework
therefore predicts that many forms of biological instability are not primarily energetic
failures but failures of admissibility compression: situations in which the effective
dimensionality of the accessible organizational space expands faster than the system's
capacity to maintain coherent constraint propagation across it.

\section{Limitations and Open Problems}
\label{sec:limitations}

The interpretive and theoretical advances in this essay carry corresponding limitations
that require explicit acknowledgment, both for intellectual honesty and to distinguish
what the framework currently establishes from what it merely suggests.

The framework is phenomenological rather than mechanistic. The RSVP equations describe
the organizational logic of constraint redistribution at an abstract level; they do not
specify the molecular or cellular mechanisms that implement the fields $\Phi$, $\mathbf{v}$,
and $S$ in any given biological substrate. This is deliberate: the framework aims to
capture organizational invariants that persist across substrates precisely because it
abstracts away from substrate-specific mechanism. But this strength is simultaneously
a limitation. The equations cannot be used to make quantitative predictions about ERK
kinetics, ZO-1 phosphorylation rates, or specific synaptic weight distributions without
additional mechanistic elaboration that the current framework does not provide.

The governing equations are underdetermined and not yet experimentally parameterized.
The coupling constants $\lambda, \alpha, \beta, \gamma, \delta, \kappa$ are introduced
as positive reals without estimated values or proposed measurement protocols. The
functional forms of $R(\Phi, S)$ and $Q(\mathbf{v})$ are left unspecified. Until these
are constrained by data, the equations function as qualitative organizational scaffolding
rather than as a quantitatively predictive model. The Signature Predictions derived in
Section~\ref{sec:predictions} are structural rather than quantitative for this reason.

Cross-scale applicability may reflect abstraction-level similarity rather than genuine
dynamical equivalence. The claim that neural developmental trajectories and epithelial
collective migration instantiate the same organizational invariants rests on structural
analogy at the level of the $(\Phi, \mathbf{v}, S)$ triple. This analogy is nontrivial
and generates testable predictions, but it does not establish that the underlying
dynamics are governed by the same equations with the same parameters. The possibility
that the shared mathematical structure is a consequence of the level of abstraction
chosen --- and would dissolve at a more detailed mechanistic level --- cannot currently
be ruled out. The invariant $\mathcal{I}$ in equation~(\ref{eq:invariant}) is offered
as a first candidate for distinguishing genuine cross-scale dynamical homology from
abstraction-level coincidence.

The accessibility field $S$ currently lacks a universally operationalized empirical
measure. The operationalization $S \approx -\log\rho$ provides a conceptually grounded
estimator, but estimating $\rho$ reliably on high-dimensional organizational manifolds
from finite biological data involves significant methodological challenges, including
manifold dimensionality estimation, boundary effects, and dependence on the dimensionality
reduction method used to reconstruct $\mathcal{M}$. The three-layer distinction
$S_{\mathrm{geom}}/S_{\mathrm{info}}/S_{\mathrm{phys}}$ clarifies the conceptual
target but does not resolve the estimation problem.

Finally, the current biological applications remain primarily interpretive rather than
predictive except for the proposed experimental signatures in Section~\ref{sec:predictions}.
The framework was introduced to interpret two already-published findings rather than to
generate those findings. Its strongest empirical claim is the collection of predictions
in Section~\ref{sec:predictions}, which have not yet been tested. The framework will
acquire a stronger scientific standing once those predictions are evaluated experimentally.

\section{Signature Predictions: Toward Empirical Discrimination}
\label{sec:predictions}

A theoretical framework earns scientific standing not only by interpreting existing
data coherently but by generating predictions that are nontrivial, structurally motivated,
and differentiable from the predictions of competing frameworks. The interpretations
developed in preceding sections are retrospective; they show that RSVP-style concepts
fit the Richier et al.\ and Hirano et al.\ findings. The present section derives
prospective predictions that would be empirically distinguishable from those expected
under standard deficit-model psychiatry, classical reaction-diffusion theory, and
active-matter descriptions of collective migration. The goal is not to prove the framework
but to expose its structure to empirical constraint.

\subsection{Longitudinal Predictions for GAD: Dynamic Variance Restructuring}

The cross-sectional PAD variance finding admits multiple interpretations. Standard
diagnostic heterogeneity models predict stable inter-individual variance arising from
fixed individual differences in GAD severity or subtype. The RSVP framework predicts
something structurally different: the variance itself should be dynamically modulated by
developmental transitions, exhibiting transient widening followed by partial reconvergence
at periods of high admissibility flux.

\begin{definition}[Variance Trajectory]
Let $\sigma^2_{PAD}(t) = \mathrm{Var}_{GAD}[\delta(t)]$ denote the population-level
PAD variance at developmental time $t$ in a longitudinally followed GAD cohort. The
RSVP framework predicts that $\sigma^2_{PAD}(t)$ is not constant but exhibits time
dependence governed by the underlying admissibility field dynamics.
\end{definition}

The first signature prediction follows directly.

\begin{proposition}[Prediction 1: Variance Transience at Developmental Transitions]
During periods of elevated developmental admissibility flux---including late adolescence,
early adulthood, and major environmental stress transitions---the PAD variance of
high-symptom GAD cohorts should transiently widen before partially reconverging, rather
than monotonically increasing or remaining stable. Formally:
\[
\frac{d}{dt}\sigma^2_{PAD}(t) \neq 0
\]
with local maxima of $\sigma^2_{PAD}$ expected to cluster near independently identified
developmental reorganization periods. A standard heterogeneity model predicts
$\frac{d}{dt}\sigma^2_{PAD}(t) \approx 0$.
\end{proposition}

This prediction is genuinely distinctive because it assigns a temporal structure to
the variance rather than treating it as a fixed population parameter. It implies that
the same cohort followed longitudinally should show rhythmic variance dynamics rather
than monotone drift, and that these dynamics should be partially predictable from
independently measured developmental markers.

A second, stronger prediction concerns trajectory substructure within
symptom-matched cohorts. The RSVP framework holds that individuals with similar symptom severity may
nevertheless occupy divergent regions of the admissibility distribution: some will be
consolidation-dominated (low $S$, high $\Phi$, positive PAD) while others will be
plasticity-dominated (high $S$, lower $\Phi$, negative PAD). Longitudinal tracking
should reveal that these individuals follow systematically different developmental
trajectories even under matched symptom burden.

\begin{proposition}[Prediction 2: Multimodal Longitudinal Substructure]
Within longitudinally followed GAD cohorts matched for baseline symptom severity,
individual PAD trajectories should exhibit multimodal substructure rather than unimodal
dispersion around a common mean trajectory. Specifically, two distinguishable subpopulations
are predicted: one following consolidation-dominated trajectories (sustained positive
PAD, decreasing developmental plasticity) and one following plasticity-dominated
trajectories (sustained negative or fluctuating PAD, elevated developmental lability).
This bimodality should be visible in latent trajectory analyses and should correlate
with independently measurable clinical markers of rigidity versus lability.
\end{proposition}

A standard deficit model predicts unimodal trajectory distributions around a slightly
deviated mean; the multimodal substructure prediction is specific to a framework in which
the admissibility field can polarize toward opposite extremes under similar perturbation
magnitudes.

\subsection{Optogenetic Predictions for ERK Wave Dynamics: Topology over Amplitude}

The Hirano et al.\ findings establish that ERK waves drive ZO-1 redistribution and
coordinated migration. The crucial RSVP-theoretic prediction concerns not whether ERK
activity matters---this is already established---but whether the spatiotemporal
\emph{topology} of ERK activation matters independently of its integrated amplitude.

\begin{proposition}[Prediction 3: Topology-Sensitive Migration Coherence]
Let $\mathcal{M}$ denote a measure of collective migration coherence (e.g., velocity
correlation length, directional order parameter, or coordinated ZO-1 redistribution
fraction). For two optogenetic stimulation protocols delivering equal integrated ERK
activation energy $E = \int_\Omega \int_0^T u(\mathbf{r}, t)\, dt\, d\mathbf{r}$:
a traveling-wave activation protocol $u_{wave}$ and a spatially uniform or localized
static activation protocol $u_{static}$, the RSVP framework predicts
\[
\mathcal{M}(u_{wave}) > \mathcal{M}(u_{static}).
\]
A reaction-diffusion or active-matter framework that treats migration as driven by
local ERK amplitude predicts $\mathcal{M}$ to depend primarily on local activation
magnitude rather than propagation topology, yielding no systematic advantage for
wave-structured over amplitude-matched static stimulation.
\end{proposition}

This prediction tests the core claim that the constraint-propagating front, rather than
the total signaling energy, is the causally relevant organizational variable. If
tissue migration coherence is substantially greater under wave-structured optogenetic
stimulation than under amplitude-matched static stimulation, this provides evidence
that propagation topology carries organizational information beyond what local ERK
amplitude encodes. If the two protocols produce similar coherence, the simpler
amplitude-based model is sufficient.

\begin{proposition}[Prediction 4: Wave-Order Hysteresis and Field Memory]
Cells previously traversed by multiple sequential ERK activation waves should exhibit
altered thresholds for ZO-1 redistribution relative to wave-naive cells under matched
instantaneous ERK amplitude. Specifically, the redistribution threshold should decrease
monotonically with prior wave exposure up to a saturation point, reflecting persistent
modification of the local admissibility structure by prior constraint-propagating fronts.
Formally, if $\theta_n$ denotes the ERK activation threshold for ZO-1 redistribution
after $n$ prior wave traversals, then $\theta_n < \theta_0$ for $n \geq 1$, with
$\theta_n \to \theta_\infty > 0$ as $n \to \infty$.
\end{proposition}

This field-memory prediction is especially important for distinguishing the RSVP
interpretation from a purely biochemical threshold account. A purely local threshold
model predicts that the redistribution response depends only on instantaneous ERK
amplitude, so that prior wave history has no effect on subsequent redistribution
thresholds once ERK amplitude returns to baseline. The RSVP prediction is that the
passage of constraint-propagating fronts leaves structural traces in the admissibility
field that persist beyond the biochemical recovery of ERK levels, producing genuine
path dependence in the tissue's organizational response.

\subsection{Cross-Prediction: Shared Variance Signature}

A cross-system prediction connects the GAD and ERK findings through the shared variance
amplification invariant identified in Section~\ref{sec:rsvp}.

\begin{proposition}[Prediction 5: Perturbation-Induced Variance Amplification]
In any coupled-field biological system in which a constraint-propagating perturbation
is applied (pharmacological, environmental, optogenetic, or developmental), the primary
observable signature of organizational impact should be an increase in the variance
of relevant state variables across system components, preceding or accompanying any
detectable shift in the population mean. This is predicted to be especially pronounced
near the boundary between organizational regimes (near the phase transition between
junction-maintenance and podosome-invasive in the tissue case; near critical-period
boundaries in the neural developmental case).
\end{proposition}

This prediction is general enough to apply across both systems and specific enough to be
empirically distinguishable from models in which perturbations produce primarily mean-field
effects. It predicts that, in drug trials, disease-state comparisons, and optogenetic
stimulation experiments alike, variance statistics should be sensitive leading indicators
of organizational impact that precede detectable mean-field changes. If consistently
confirmed, this would constitute the strongest empirical support for treating variance
geometry as a primary organizational variable rather than as noise to be averaged away.

The experimental addendum proposed here does not exhaust the predictive scope of the RSVP
framework. It identifies the predictions most directly motivated by the two empirical
papers discussed and most sharply distinguished from standard competing accounts.
Whether these predictions are confirmed or refuted, their articulation serves the function
of making the framework scientifically negotiable: the conceptual commitments of the
essay become hypotheses about how biological systems actually behave, rather than
interpretive lenses that can accommodate any outcome.

\begin{thebibliography}{99}
\sloppy

\bibitem{Wang2026}
S. Wang, K. Lou, C. Wei et al.,
``A geometry-aware framework enhances noninvasive mapping of whole human brain dynamics,''
\textit{Nature Biomedical Engineering}, 2026.
doi:10.1038/s41551-026-01664-0.

\bibitem{Richier2026}
C. Richier, A. Zugman, A. Harrewijn et al.,
``Brain age prediction in generalized anxiety disorder using a convolutional neural network,''
\textit{Translational Psychiatry}, 2026.
doi:10.1038/s41398-026-04078-3.

\bibitem{Hirano2026}
S. Hirano, Y. Kondo, A. Kitajima et al.,
``ZO-1 shuttles between apical junctional complexes and podosomes by riding ERK activation waves,''
\textit{Nature Communications}, 2026.
doi:10.1038/s41467-026-72840-8.

\bibitem{Friston2010}
K. Friston,
``The free-energy principle: a unified brain theory?''
\textit{Nature Reviews Neuroscience},
vol.\ 11, no.\ 2, pp.\ 127--138, 2010.

\bibitem{Kelso1995}
J. A. S. Kelso,
\textit{Dynamic Patterns: The Self-Organization of Brain and Behavior}.
Cambridge, MA: MIT Press, 1995.

\bibitem{Nicolis1977}
G. Nicolis and I. Prigogine,
\textit{Self-Organization in Nonequilibrium Systems}.
New York: Wiley, 1977.

\bibitem{Kauffman1993}
S. A. Kauffman,
\textit{The Origins of Order: Self-Organization and Selection in Evolution}.
New York: Oxford University Press, 1993.

\bibitem{Turing1952}
A. M. Turing,
``The chemical basis of morphogenesis,''
\textit{Philosophical Transactions of the Royal Society B},
vol.\ 237, no.\ 641, pp.\ 37--72, 1952.

\bibitem{Murray2003}
J. D. Murray,
\textit{Mathematical Biology II: Spatial Models and Biomedical Applications},
3rd ed.
New York: Springer, 2003.

\bibitem{Bak1987}
P. Bak, C. Tang, and K. Wiesenfeld,
``Self-organized criticality: an explanation of $1/f$ noise,''
\textit{Physical Review Letters},
vol.\ 59, no.\ 4, pp.\ 381--384, 1987.

\bibitem{Tononi2004}
G. Tononi,
``An information integration theory of consciousness,''
\textit{BMC Neuroscience},
vol.\ 5, no.\ 42, 2004.

\bibitem{Barabasi2016}
A.-L. Barab{\'a}si,
\textit{Network Science}.
Cambridge: Cambridge University Press, 2016.

\bibitem{Raichle2015}
M. E. Raichle,
``The brain's default mode network,''
\textit{Annual Review of Neuroscience},
vol.\ 38, pp.\ 433--447, 2015.

\bibitem{Levin2012}
M. Levin,
``Morphogenetic fields in embryogenesis, regeneration, and cancer:
non-local control of complex patterning,''
\textit{BioSystems},
vol.\ 109, no.\ 3, pp.\ 243--261, 2012.

\bibitem{Camazine2003}
S. Camazine et al.,
\textit{Self-Organization in Biological Systems}.
Princeton: Princeton University Press, 2003.

\bibitem{Haken1983}
H. Haken,
\textit{Synergetics: An Introduction},
3rd ed.
Berlin: Springer, 1983.

\bibitem{Maturana1980}
H. R. Maturana and F. J. Varela,
\textit{Autopoiesis and Cognition: The Realization of the Living}.
Dordrecht: Reidel, 1980.

\bibitem{DeLanda2002}
M. DeLanda,
\textit{Intensive Science and Virtual Philosophy}.
London: Continuum, 2002.

\bibitem{Thompson2007}
E. Thompson,
\textit{Mind in Life: Biology, Phenomenology, and the Sciences of Mind}.
Cambridge, MA: Harvard University Press, 2007.

\bibitem{Engel2001}
A. K. Engel, P. Fries, and W. Singer,
``Dynamic predictions: oscillations and synchrony in top-down processing,''
\textit{Nature Reviews Neuroscience},
vol.\ 2, no.\ 10, pp.\ 704--716, 2001.

\bibitem{Singer1999}
W. Singer,
``Neuronal synchrony: a versatile code for the definition of relations?''
\textit{Neuron},
vol.\ 24, no.\ 1, pp.\ 49--65, 1999.

\bibitem{Meinhardt1982}
H. Meinhardt,
\textit{Models of Biological Pattern Formation}.
London: Academic Press, 1982.

\bibitem{Smoller2020}
J. W. Smoller,
``The use of electronic health records for psychiatric phenotyping and genomics,''
\textit{American Journal of Medical Genetics Part B},
vol.\ 183, no.\ 1, pp.\ 3--12, 2020.

\bibitem{Goodwin1994}
B. Goodwin,
\textit{How the Leopard Changed Its Spots}.
New York: Scribner, 1994.

\bibitem{Gimperlein2026}
H. Gimperlein, M. Grinfeld, R. J. Knops, and M. Slemrod,
``On action rate admissibility criteria,''
\textit{Zeitschrift f\"ur angewandte Mathematik und Physik},
vol.\ 77, article 57, 2026.
doi:10.1007/s00033-025-02705-5.

\bibitem{Chemnitz2025}
R. Chemnitz and D. Dragi\v{c}evi\'c,
``Characterizing nonuniform hyperbolicity by Mather-type admissibility,''
\textit{Journal of Dynamics and Differential Equations},
vol.\ 37, pp.\ 3197--3216, 2025.
doi:10.1007/s10884-024-10398-z.

\bibitem{Barreira2024}
L. Barreira and C. Valls,
``Admissibility via induced delay equations,''
\textit{Qualitative Theory of Dynamical Systems},
vol.\ 23, article 229, 2024.
doi:10.1007/s12346-024-01086-w.

\bibitem{Choudhury2025}
B. S. Choudhury, N. Metiya, A. Kundu, and S. Kundu,
``Multivalued coupled coincidence point results using admissibility,''
\textit{Vestnik St.\ Petersburg University: Mathematics},
vol.\ 58, no.\ 3, pp.\ 408--418, 2025.
doi:10.1134/S1063454125700402.

\bibitem{Topey2026}
B. Topey,
``Whence admissibility constraints? From inferentialism to tolerance,''
\textit{Journal of Philosophical Logic},
vol.\ 55, pp.\ 411--435, 2026.
doi:10.1007/s10992-026-09836-8.

\bibitem{ElAbbassi2025}
M. El Abbassi and C. Ziti,
``Comparison of a hyperbolic equation and system:
admissibility and non-classical waves,''
\textit{International Journal of Applied and Computational Mathematics},
vol.\ 11, article 140, 2025.
doi:10.1007/s40819-025-01936-4.

\bibitem{Monjon2023}
O. Monjon, J. Scherer, and F. Sterck,
``Admissibility of localizations of crossed modules,''
\textit{Applied Categorical Structures},
vol.\ 31, article 37, 2023.
doi:10.1007/s10485-023-09738-9.

\end{thebibliography}

\end{document}
